% Shared packages and notation. The official ICLR style is not modified.
\usepackage{amsmath,amssymb,mathrsfs}
\usepackage{microtype}
\usepackage{booktabs,array,tabularx}
\usepackage{graphicx,xcolor}
\usepackage{colortbl}
% Near-white washes; typography, rather than saturated fills, carries emphasis.
\definecolor{mainProp}{HTML}{F8FAFC}
\definecolor{mainGeneral}{HTML}{F8FBFA}
\definecolor{mainMedical}{HTML}{FDFAF8}
\definecolor{mainOurs}{HTML}{F0F4FA}
% User-selected VQA main-table palette; 3D tables retain their existing colors.
\definecolor{vqaProp}{HTML}{F6D6D6}
\definecolor{vqaGeneral}{HTML}{C0E0F8}
\definecolor{vqaMedical}{HTML}{C2E2C2}
% Table 1: retain the selected hues at 25% color on a white background.
\colorlet{vqaProp}{vqaProp!25!white}
\colorlet{vqaGeneral}{vqaGeneral!25!white}
\colorlet{vqaMedical}{vqaMedical!25!white}
% Shared by the English and Chinese manuscripts.
% Columns: PMC-VQA, VQA-RAD, SLAKE, PathVQA. Reports are in figures/report_ratescore.
% VQA: user-supplied comparison image, except Hulu-Med (Jiang et al., Table 1).
% Report values are preserved in figures/report_results.json.
% CoEvoMed-8B: user-confirmed same-checkpoint values; IU-Xray=68.22; PathVQA author-corrected to 69.40 on 2026-09-16.
% CoEvoMed-32B: six author-supplied scores on 2026-09-13; VQA-RAD corrected by author to 82.71 on 2026-09-17; report lexical metrics pending.
% A dash is unavailable, not zero.
% Qwen3.5-9B: author-supplied overall accuracy; figures/qwen35_9b_vqa_results.json.
% Qwen3-VL-8B: former Qwen3 (base) row, identity confirmed by author 2026-09-18.
% Qwen3-VL-32B: Qwen3-VL-32B-Instruct, author-supplied full metrics on 2026-09-18; 2D_result/other_model.json.
% MedVLSynther-7B: Huang et al., ICLR 2026, Table 6 (published results).
\newcommand{\MainResultsRows}{%
\midrule
\multicolumn{5}{c}{\itshape Proprietary Models}\\
\midrule
\rowcolor{vqaProp} Gemini 3 Flash & 62.62 & 69.62 & 74.55 & 52.03\\
\rowcolor{vqaProp} GPT-5.6-sol & 58.89 & 67.18 & 69.15 & 44.72\\
\rowcolor{vqaProp} Kimi-K3 & 64.48 & 73.61 & 77.13 & 53.90\\
\rowcolor{vqaProp} MiniMax-M3 & 57.33 & 68.07 & 71.44 & 57.80\\
\rowcolor{vqaProp} Qwen3.8-Max & 64.58 & 76.81 & 85.05 & 64.11\\
\midrule
\multicolumn{5}{c}{\itshape General-purpose Multimodal VLMs}\\
\midrule
\rowcolor{vqaGeneral} Qwen3-VL-8B & 54.13 & 56.76 & 64.84 & 42.77\\
\rowcolor{vqaGeneral} Qwen3.5-9B & 60.71 & 68.74 & 71.87 & 49.64\\
\rowcolor{vqaGeneral} Llama3.2-11B & 48.10 & 58.80 & 65.80 & 32.90\\
\rowcolor{vqaGeneral} InternVL3-14B & 54.10 & 66.30 & 72.80 & 48.00\\
\rowcolor{vqaGeneral} Qwen3.6-27B & 62.50 & 71.62 & 75.07 & 46.93\\
\rowcolor{vqaGeneral} Qwen3.8-27B & 60.33 & 68.51 & 73.93 & 45.11\\
\rowcolor{vqaGeneral} Qwen3-VL-32B & 57.24 & 65.85 & 71.87 & 44.28\\
\rowcolor{vqaGeneral} InternVL3-38B & 56.60 & 65.40 & 72.70 & 51.00\\
\midrule
\multicolumn{5}{c}{\itshape Medical Multimodal VLMs}\\
\midrule
\multicolumn{5}{l}{\itshape Models $<10$B}\\
\rowcolor{vqaMedical} MedGemma-4B-IT & 49.20 & 72.30 & 78.20 & 48.10\\
\rowcolor{vqaMedical} HuatuoGPT-V-7B & 53.10 & 67.60 & 68.10 & 44.80\\
\rowcolor{vqaMedical} Lingshu-7B & 56.30 & 67.90 & 83.10 & 61.90\\
\rowcolor{vqaMedical} GMAI-VL-7B & 56.09 & 57.21 & 59.69 & 39.05\\
\rowcolor{vqaMedical} MedVLSynther-7B & 55.43 & 77.57 & 72.36 & 65.56\\
\rowcolor{vqaMedical} Hulu-Med-7B & 66.80 & 78.00 & 86.80 & 65.60\\
\rowcolor{vqaMedical} \textbf{CoEvoMed (8B)} & \textbf{67.95} & \textbf{82.48} & \textbf{91.60} & \textbf{69.40}\\
\multicolumn{5}{l}{\itshape Models $>10$B}\\
\rowcolor{vqaMedical} HuatuoGPT-V-34B & 56.60 & 61.40 & 69.50 & 44.40\\
\rowcolor{vqaMedical} Lingshu-32B & 57.90 & 76.70 & 86.70 & 65.50\\
\rowcolor{vqaMedical} MedDr-40B & 13.90 & 65.20 & 66.40 & 53.50\\
\rowcolor{vqaMedical} Hulu-Med-32B & 69.40 & 81.40 & 85.70 & 67.30\\
\rowcolor{vqaMedical} \textbf{CoEvoMed (32B)} & \textbf{70.20} & \textbf{82.71} & \textbf{92.60} & \textbf{69.60}\\
\bottomrule
}

% Shared source-tagged 3D comparisons; scores scaled by 100.
\newcommand{\MThreeRows}{%
\rowcolor{mainGeneral} RadFM\textsuperscript{*} & 19.60 & 28.70 & 16.80 & 18.90 & 14.90\\
\rowcolor{mainGeneral} Qwen2.5VL-72B\textsuperscript{*} & 94.80 & 50.00 & 62.40 & 62.50 & 59.80\\
\rowcolor{mainGeneral} Lingshu-32B\textsuperscript{*} & 93.80 & 61.60 & 60.90 & 61.50 & 57.60\\
\rowcolor{mainGeneral} M3D-Llama2-7B\textsuperscript{*} & 98.80 & 79.80 & 74.80 & 66.70 & 58.90\\
\rowcolor{mainGeneral} Hulu-Med-7B\textsuperscript{*} & 99.00 & 88.00 & 78.90 & 70.70 & 67.30\\
\rowcolor{mainGeneral} Qwen3.5-9B & 96.50 & 67.25 & 59.50 & 59.06 & 57.33\\
\rowcolor{mainGeneral} Qwen3.6-27B & 96.45 & 69.95 & 63.75 & 60.84 & 59.56\\
\rowcolor{mainGeneral} Qwen3.8-27B & 95.60 & 36.55 & 63.70 & 57.92 & 56.39\\
\rowcolor{mainGeneral} Med3DVLM-7B & 99.25 & 90.65 & 80.00 & 73.28 & 66.96\\
\rowcolor{mainGeneral} Photon-3B-S1 & 51.50 & 26.50 & 42.80 & 48.73 & 48.65\\
\rowcolor{mainOurs} \textbf{CoEvoMed} & \textbf{99.50} & \textbf{91.90} & \textbf{82.20} & \textbf{73.88} & \textbf{67.80}\\
}
\newcommand{\RadThreeRows}{%
\rowcolor{mainGeneral} Qwen2.5VL-72B\textsuperscript{*} & 15.50 & 25.20 & 16.90 & 27.30 & 6.70 & 11.60\\
\rowcolor{mainGeneral} M3D (Phi3-4B)\textsuperscript{*} & 16.30 & 23.20 & 15.10 & 23.20 & 2.50 & 5.60\\
\rowcolor{mainGeneral} Lingshu-7B\textsuperscript{*} & 9.30 & 21.60 & 10.40 & 22.20 & 0.90 & 8.40\\
\rowcolor{mainGeneral} M3D (Llama2-7B)\textsuperscript{*} & 10.70 & 20.80 & 9.10 & 18.60 & 15.90 & 23.20\\
\rowcolor{mainGeneral} Lingshu-32B\textsuperscript{*} & 7.60 & 16.10 & 7.80 & 16.30 & 0.90 & 4.80\\
\rowcolor{mainGeneral} Hulu-Med-7B\textsuperscript{*} & 23.40 & 35.90 & 23.00 & 37.10 & 23.50 & 26.00\\
\rowcolor{mainGeneral} Qwen3.5-9B & 12.85 & 25.04 & 14.07 & 24.18 & 7.05 & 15.17\\
\rowcolor{mainGeneral} Qwen3.6-27B & 22.47 & 29.68 & 23.09 & 30.76 & 2.89 & 7.99\\
\rowcolor{mainGeneral} Qwen3.8-27B & 12.70 & 26.63 & 14.62 & 30.00 & 13.08 & 13.77\\
\rowcolor{mainGeneral} Med3DVLM-7B & 0.68 & 6.23 & 0.77 & 5.79 & 0.11 & 0.46\\
\rowcolor{mainGeneral} Photon-3B-S2 (local) & 37.93 & 37.58 & 48.09 & 47.26 & 29.22 & 31.13\\
\rowcolor{mainOurs} \textbf{CoEvoMed} & \textbf{38.60} & \textbf{44.86} & \textbf{52.36} & \textbf{57.06} & \textbf{37.83} & \textbf{39.30}\\
}
\newcommand{\AmossRows}{%
\rowcolor{mainGeneral} Qwen3.6-27B & 46.44 & 15.49 & 18.88 & 13.16\\
\rowcolor{mainGeneral} Hulu-Med-7B & 64.35 & 35.84 & 37.49 & 30.51\\
\rowcolor{mainGeneral} M3D-Llama3.1-8B\textsuperscript{*} & 57.00 & 26.40 & 37.90 & \textbf{35.50}\\
\rowcolor{mainGeneral} M3D-Mistralv0.3-7B\textsuperscript{*} & 56.90 & 25.80 & 37.40 & 34.60\\
\rowcolor{mainGeneral} Qwen2.5VL-32B\textsuperscript{*} & 52.90 & 24.60 & 19.10 & 21.10\\
\rowcolor{mainGeneral} Qwen3.5-9B & 53.68 & 29.26 & 21.40 & 20.04\\
\rowcolor{mainGeneral} Qwen3.8-27B & 50.02 & 23.06 & 18.93 & 16.34\\
\rowcolor{mainGeneral} Med3DVLM-7B & 42.84 & 11.63 & 15.66 & 10.34\\
\rowcolor{mainGeneral} Lingshu-7B & 50.75 & 24.43 & 19.89 & 17.75\\
\rowcolor{mainGeneral} Lingshu-32B & 45.22 & 24.45 & 17.43 & 16.75\\
\rowcolor{mainGeneral} M3D-LaMed-4B & 39.00 & 4.32 & 12.57 & 5.87\\
\rowcolor{mainOurs} \textbf{CoEvoMed} & \textbf{66.97} & \textbf{39.30} & \textbf{40.40} & 31.91\\
}

\usepackage{tikz}
\usetikzlibrary{arrows.meta,positioning,fit,backgrounds}
\usepackage{placeins}
\usepackage{float}
% Pseudocode uses the existing float/tabularx packages; captions remain below.
\floatstyle{plain}
\newfloat{algorithm}{htbp}{loa}
\floatname{algorithm}{Algorithm}
\usepackage{wrapfig}
% Keep natural vertical spacing when an indivisible table moves to the next page.
% Do not stretch headings and paragraph gaps to force page-bottom alignment.
\AtBeginDocument{\raggedbottom}
\usepackage{hyperref}
% Allow long reference URLs to wrap without stretching the surrounding words.
\usepackage{xurl}
\hypersetup{colorlinks=true,linkcolor=blue!45!black,citecolor=blue!45!black,urlcolor=blue!45!black,pdfauthor={Anonymous authors}}
\newcommand{\tbd}{\textcolor{black!65}{\textsf{TBD}}}
\newcommand{\obs}{o}
\newcommand{\hist}{h}
\newcommand{\strategies}{\Sigma}
\newcommand{\evol}{\mathscr E}
\newcommand{\cellset}{\mathcal C}
\newcommand{\source}{\mathcal S}
\newcommand{\trainset}{\mathcal D}
\newcommand{\E}{\mathbb E}
\newcommand{\ind}{\mathbf 1}
\newcolumntype{L}[1]{>{\raggedright\arraybackslash}p{#1}}
\newcolumntype{Y}{>{\raggedright\arraybackslash}X}

% Language-specific labels are supplied by each main file.
\newcommand{\frameworkfigure}{%
\begin{tikzpicture}[
  box/.style={draw=black!60,rounded corners=2pt,align=center,font=\small,
    minimum height=0.80cm,text width=3.45cm,inner sep=4pt},
  arrow/.style={-{Latex[length=2mm]},thick,draw=black!70},
  feed/.style={arrow,dashed},node distance=0.7cm and 0.6cm]
  \node[box,fill=blue!6] (feedback) {\FigFeedback};
  \node[box,fill=blue!6,right=of feedback] (policy) {\FigPolicy};
  \node[box,fill=orange!9,right=of policy] (synth) {\FigSynthesis};
  \node[box,fill=orange!9,below=of synth] (student) {\FigStudent};
  \node[box,fill=orange!9,left=of student] (strategy) {\FigStrategy};
  \node[box,fill=black!3,left=of strategy] (record) {\FigRecord};
  \draw[arrow] (feedback) -- (policy);
  \draw[arrow] (policy) -- (synth);
  \draw[arrow] (synth) -- (student);
  \draw[arrow] (student.south) -- ++(0,-0.45) -| (record.south);
  \draw[arrow] (record) -- (feedback);
  \draw[feed] (feedback.south) -- (strategy.north west);
  \draw[arrow] (strategy.north east) -- (synth.south west);
\end{tikzpicture}%
}

% An honest placeholder for future measurements, not an illustrative result curve.
\newcommand{\trajectoryplaceholder}{%
\begin{tikzpicture}[font=\small]
  \draw[draw=black!45,rounded corners=2pt] (0,2.5) rectangle (12.8,4.6);
  \node at (6.4,4.1) {\FigCapability};
  \node[align=center,text width=11.5cm,text=black!60] at (6.4,3.25) {\FigPending};
  \draw[draw=black!45,rounded corners=2pt] (0,0) rectangle (6.2,2.2);
  \draw[draw=black!45,rounded corners=2pt] (6.6,0) rectangle (12.8,2.2);
  \node at (3.1,1.7) {\FigVQA};
  \node at (9.7,1.7) {\FigReport};
  \node[align=center,text width=5.5cm,text=black!60] at (3.1,0.9) {\FigPending};
  \node[align=center,text width=5.5cm,text=black!60] at (9.7,0.9) {\FigPending};
\end{tikzpicture}%
}
